\documentclass[12pt,a4paper]{article}
\usepackage[nowatermark]{mathnexus}
\usepackage{enumitem}
\usepackage{pgfplots}\pgfplotsset{compat=1.17}
\mnsetup{type={Série d'exercices},niveau={Terminale S2},matiere={Physique},chapnum={P2},chapitre={Bases de la dynamique},auteur={MathNexus},annee={2026--2027}}
\newgeometry{left=1.4cm,right=1.4cm,top=3.1cm,bottom=1.9cm,headheight=2.0cm,headsep=0.35cm,footskip=0.8cm}
\fancyhead[L]{\raisebox{-0.35\height}{\includegraphics[height=1.75cm]{logo.png}}}
\fancyhead[C]{}
\fancyhead[R]{\sffamily\color{mnnavy}\begin{tabular}{@{}r@{}}{\large\bfseries Terminale S2 \,$\cdot$\, Physique}\\[1pt]
  {\color{mnblue}Série P2 : Bases de la dynamique}\\[1pt]
  {\small\color{mngrey}Année scolaire 2026--2027}\end{tabular}}
\newcommand{\vect}[1]{\overrightarrow{#1}}
\newcommand{\U}[1]{\,\mathrm{#1}}
\providecommand{\ms}{\,\mathrm{m\cdot s^{-1}}}
\providecommand{\mss}{\,\mathrm{m\cdot s^{-2}}}
\providecommand{\dd}{\mathrm{d}}
\newcommand{\Exo}[3]{\par\Needspace{6\baselineskip}\medskip\noindent{\sffamily\bfseries\color{mnblue}Exercice #1 \,·\, #2}\hfill{\sffamily\small\color{mnnavy}\ifx\relax#3\relax\else #3 pts\fi}\par\nointerlineskip{\color{mnblue}\noindent\rule{\linewidth}{0.5pt}}\par\smallskip}
\setlength{\parskip}{3pt}
\begin{document}
\fontsize{12.5}{15.2}\selectfont
\mntitreserie{P2}{Bases de la dynamique}

\sujet{Partie A · Centre d'inertie et référentiels}

\Exo{1}{Le système Terre-Lune}{}
La masse de la Terre vaut 81 fois celle de la Lune ; la distance entre leurs centres vaut $384\,000\U{km}$. Où est le centre d'inertie du système Terre-Lune ? Est-il à l'intérieur de la Terre (rayon $6400\U{km}$) ?


\Exo{2}{Comparer des quantités de mouvement}{}
\begin{enumerate}[label=\textbf{\arabic*.},leftmargin=*,start=1]\item Calculer la quantité de mouvement d'une voiture de $1000\U{kg}$ roulant à $72\U{km/h}$.
\item À quelle vitesse un camion de $10\U{t}$ a-t-il la même quantité de mouvement ?
\end{enumerate}

\Exo{3}{Choisir le référentiel}{}
Quel référentiel choisir pour étudier : a) un satellite autour de la Terre ; b) la Terre autour du Soleil ; c) un ballon de football ?


\Exo{4}{L'ascenseur}{}
Un ascenseur de $400\U{kg}$ monte à vitesse constante. On prend $g=10\mss$. Calculer la tension du câble.


\Exo{5}{L'haltère (synthèse)}{10}
Un haltère est formé d'une barre homogène de masse $1\U{kg}$ et de deux boules : $m_1=2\U{kg}$ et $m_2=3\U{kg}$, dont les centres sont aux extrémités de la barre, distants de $50\U{cm}$.


\begin{enumerate}[label=\textbf{\arabic*.},leftmargin=*,start=1]\item Déterminer la position du centre d'inertie $G$ par rapport au centre de la boule 1.
\item L'haltère est déplacé en translation à $0{,}5\ms$. Calculer sa quantité de mouvement.
\item Un élève affirme : « $G$ est au milieu de la barre ». Expliquer son erreur.
\end{enumerate}

\Exo{6}{Le bloc tiré à vitesse constante (synthèse)}{10}
Un bloc de masse $4\U{kg}$ est tiré sur un plan horizontal par une force horizontale $\vec F$ de valeur $12\U{N}$. Il se déplace en ligne droite à vitesse constante. On prend $g=10\mss$.


\begin{enumerate}[label=\textbf{\arabic*.},leftmargin=*,start=1]\item Énoncer le principe de l'inertie.
\item Faire le bilan des forces. En déduire la force de frottement $\vec f$ et la composante normale $\vec R_N$ de la réaction.
\item Calculer la valeur de la réaction totale $\vec R=\vec R_N+\vec f$ et l'angle qu'elle fait avec la verticale.
\end{enumerate}

\sujet{Partie B · Le théorème du centre d'inertie}

\Exo{7}{L'ascenseur qui démarre}{}
Un ascenseur de $500\U{kg}$ démarre vers le haut avec une accélération de $2\mss$. On prend $g=10\mss$. Calculer la tension du câble. La comparer au poids.


\Exo{8}{Sans frottement}{}
Un solide de $5\U{kg}$, au repos sur un plan horizontal sans frottement, est poussé par une force horizontale constante de $30\U{N}$. Calculer son accélération, puis sa vitesse et la distance parcourue après $2\U{s}$.


\Exo{9}{La montée}{}
Un solide est lancé vers le haut d'un plan incliné de $30^{\circ}$ (sans frottement) avec une vitesse de $6\ms$. On prend $g=10\mss$. Quelle distance parcourt-il avant de s'arrêter, et en combien de temps ?


\Exo{10}{Mesurer un frottement}{}
Un solide de $0{,}5\U{kg}$ glisse sur un plan incliné de $20^{\circ}$. On mesure $a=2{,}4\mss$. On prend $g=10\mss$. Calculer la force de frottement.


\Exo{11}{Le plan incliné puis le plan horizontal (synthèse)}{10}
Un solide de masse $m=1\U{kg}$ part sans vitesse du sommet $A$ d'un plan incliné de $30^{\circ}$, de longueur $AB=4\U{m}$. Sur tout le trajet, une force de frottement constante $f=2\U{N}$ s'oppose au mouvement. On prend $g=10\mss$.


\begin{enumerate}[label=\textbf{\arabic*.},leftmargin=*,start=1]\item Appliquer le TCI sur $AB$ et calculer l'accélération $a_1$ ainsi que la réaction normale $R_N$.
\item Calculer la vitesse en $B$ et la durée du trajet $AB$.
\item Le solide aborde en $B$ un plan horizontal, avec le même frottement. Calculer sa nouvelle accélération et la distance parcourue avant l'arrêt.
\end{enumerate}

\Exo{12}{La luge tirée par une corde (synthèse)}{10}
Une luge de $20\U{kg}$ est tirée sur un sol horizontal par une corde inclinée de $30^{\circ}$ au-dessus de l'horizontale. La tension vaut $T=100\U{N}$ et les frottements $f=50\U{N}$. On prend $g=10\mss$.


\begin{enumerate}[label=\textbf{\arabic*.},leftmargin=*,start=1]\item Faire le bilan des forces et écrire le TCI.
\item Projeter sur l'horizontale : calculer l'accélération.
\item Projeter sur la verticale : calculer la réaction normale du sol. Pourquoi est-elle inférieure au poids ?
\item Partie du repos, quelle vitesse la luge atteint-elle en $5\U{s}$ ?
\end{enumerate}

\sujet{Partie C · Les théorèmes énergétiques}

\Exo{13}{La distance de freinage}{}
Une voiture de $1000\U{kg}$ roulant à $20\ms$ freine jusqu'à l'arrêt. La force de freinage, constante, vaut $8000\U{N}$.


\begin{enumerate}[label=\textbf{\arabic*.},leftmargin=*,start=1]\item Calculer le travail de la force de freinage.
\item En déduire la distance de freinage.
\end{enumerate}

\Exo{14}{La chute libre}{}
Une pierre est lâchée sans vitesse d'une hauteur de $20\U{m}$. On néglige l'air. On prend $g=10\mss$. Avec quelle vitesse touche-t-elle le sol ?


\Exo{15}{Le pendule simple}{}
Un pendule de longueur $\ell=1\U{m}$ est lâché sans vitesse quand le fil fait $60^{\circ}$ avec la verticale. On néglige les frottements. On prend $g=10\mss$. Calculer la vitesse de la masse au passage par la verticale.


\Exo{16}{Le toboggan}{}
Un objet de $0{,}5\U{kg}$ part sans vitesse d'une hauteur de $1{,}8\U{m}$ sur un toboggan et arrive en bas à $5\ms$. On prend $g=10\mss$. Y a-t-il des frottements ? Calculer l'énergie perdue.


\Exo{17}{Le skieur (synthèse)}{10}
Un skieur de $80\U{kg}$ part sans vitesse du sommet $A$ d'une piste rectiligne $AB$ de longueur $50\U{m}$ ; la dénivellation entre $A$ et $B$ vaut $10\U{m}$. Sur $AB$, les frottements sont négligeables. Il continue ensuite sur une piste horizontale $BC$ où une force de frottement constante de $200\U{N}$ s'oppose au mouvement. On prend $g=10\mss$.


\begin{enumerate}[label=\textbf{\arabic*.},leftmargin=*,start=1]\item Calculer le travail du poids entre $A$ et $B$. Que vaut le travail de la réaction ?
\item Calculer la vitesse en $B$ avec le TEC.
\item Quelle distance le skieur parcourt-il sur $BC$ avant de s'arrêter ?
\end{enumerate}

\Exo{18}{La balle et l'énergie mécanique (synthèse)}{10}
Une balle de $200\U{g}$ est lancée verticalement vers le haut à $8\ms$ depuis une hauteur de $1{,}5\U{m}$. On prend $E_p=0$ au sol. On prend $g=10\mss$.


\begin{enumerate}[label=\textbf{\arabic*.},leftmargin=*,start=1]\item Calculer l'énergie mécanique au départ.
\item Sans frottement, calculer la hauteur maximale atteinte et la vitesse en touchant le sol.
\item En réalité, la balle monte jusqu'à $4{,}2\U{m}$. Calculer l'énergie perdue pendant la montée et la valeur moyenne de la force de frottement de l'air.
\end{enumerate}

\sujet{Partie D · La conservation de la quantité de mouvement}

\Exo{19}{Le canon}{}
Un canon de $500\U{kg}$ tire horizontalement un obus de $5\U{kg}$ à $200\ms$. Calculer la vitesse de recul du canon.


\Exo{20}{Sauter d'une barque}{}
Un homme de $70\U{kg}$ saute horizontalement à $2\ms$ d'une barque de $140\U{kg}$ immobile. Avec quelle vitesse la barque part-elle ? (On néglige les frottements de l'eau.)


\Exo{21}{Masses égales}{}
Deux boules de même masse ; la première arrive à $v_1$ sur la seconde, immobile. Le choc est élastique. Que se passe-t-il ?


\Exo{22}{La balle dans le bloc}{}
Une balle de $20\U{g}$ à $300\ms$ s'enfonce dans un bloc de bois de $980\U{g}$ immobile sur la glace.


\begin{enumerate}[label=\textbf{\arabic*.},leftmargin=*,start=1]\item Calculer la vitesse de l'ensemble juste après le choc.
\item Quel pourcentage de l'énergie cinétique a été perdu ?
\end{enumerate}

\Exo{23}{Le patineur et le ballon (synthèse)}{10}
Un patineur de $50\U{kg}$, immobile sur la glace, lance horizontalement un ballon de $2\U{kg}$ à $10\ms$. On prend $g=10\mss$.


\begin{enumerate}[label=\textbf{\arabic*.},leftmargin=*,start=1]\item Montrer que le système {patineur + ballon} est pseudo-isolé pendant le lancer.
\item Calculer la vitesse du patineur juste après le lancer.
\item Le patineur s'arrête après avoir glissé sur $2\U{m}$. En supposant la force de frottement constante, calculer sa valeur (par le TEC).
\end{enumerate}

\Exo{24}{Les deux pendules (synthèse)}{10}
Deux boules identiques de $100\U{g}$ sont suspendues à des fils de même longueur et se touchent au repos. On écarte la boule $B_1$ et on la lâche sans vitesse d'une hauteur $h=20\U{cm}$ au-dessus de sa position de repos. On néglige les frottements. On prend $g=10\mss$.


\begin{enumerate}[label=\textbf{\arabic*.},leftmargin=*,start=1]\item Calculer la vitesse $v_1$ de $B_1$ juste avant le choc.
\item Le choc est élastique. Donner les vitesses juste après le choc. À quelle hauteur remonte $B_2$ ?
\item On recommence avec des boules en pâte à modeler qui restent collées (choc mou). Calculer leur vitesse commune, la hauteur atteinte et l'énergie perdue dans le choc.
\end{enumerate}

\newpage\sujet{Évaluation 1 du chapitre (120 min)}

\Exo{25}{La piste ABC (questions 1 à 3 : bac 2024)}{7}
\emph{Les questions 1 à 3 sont extraites du bac S2 2024 (exercice 3, partie 3.1).} Un objet se déplace sur une piste située dans un plan vertical : $AB$ est rectiligne, de longueur $\ell=1\U{m}$, inclinée d'un angle $\alpha=30^{\circ}$ par rapport à l'horizontale ; $BC$ est rectiligne, horizontale, de longueur $\ell=1\U{m}$. Sur $AB$ et $BC$ s'exercent des forces de frottement équivalentes à une force unique $\vec f$ parallèle à la piste, opposée au vecteur vitesse et d'intensité constante. Un solide $S$, ponctuel, de masse $m=200\U{g}$, part du point $A$ sans vitesse à la date $t=0$. On prendra $g=10\U{SI}$.


\begin{enumerate}[label=\textbf{\arabic*.},leftmargin=*,start=1]\item Exprimer la norme $v_B$ de la vitesse du solide en $B$ en fonction de $m$, $g$, $\ell$, $f$ et $\alpha$.
\item Calculer l'intensité de $\vec f$ pour que le solide arrive en $C$ avec une vitesse nulle.
\item Déterminer la date à laquelle le solide arrive en $B$.
\item Par le TCI, calculer l'accélération sur $BC$ et la durée du trajet $BC$.
\item Calculer la composante normale $\vec R_N$ de la réaction sur $AB$, puis la valeur de la réaction totale.
\item On prend $E_p=0$ au niveau de $BC$. Calculer l'énergie mécanique en $A$ et en $C$. Comparer la perte au travail des frottements.
\end{enumerate}

\Exo{26}{Le pendule balistique}{6}
Un fusil de $2\U{kg}$ tire horizontalement une balle de $10\U{g}$, qui s'enfonce dans un bloc de bois de $990\U{g}$ suspendu à un fil. Juste après le choc, le bloc (avec la balle) remonte d'une hauteur $h=5\U{cm}$. On néglige les frottements. On prend $g=10\mss$.


\begin{enumerate}[label=\textbf{\arabic*.},leftmargin=*,start=1]\item Calculer la vitesse $v'$ du bloc juste après le choc.
\item Quel est le type de choc ? En déduire la vitesse $v$ de la balle avant le choc.
\item Calculer l'énergie cinétique perdue dans le choc et la comparer à l'énergie initiale.
\item Calculer la vitesse de recul du fusil.
\item Calculer la vitesse du centre d'inertie du système {balle + bloc} juste avant et juste après le choc. Commenter.
\end{enumerate}

\Exo{27}{Le téléski}{7}
Un skieur de $70\U{kg}$ est tiré par la perche d'un téléski, parallèle à une piste inclinée de $20^{\circ}$ sur l'horizontale. Les frottements équivalent à une force constante de $50\U{N}$. On prend $g=10\mss$. ($\sin20^{\circ}\approx0{,}342$ ; $\cos20^{\circ}\approx0{,}940$)


\begin{enumerate}[label=\textbf{\arabic*.},leftmargin=*,start=1]\item Le skieur monte à vitesse constante. Énoncer le principe utilisé, puis calculer la tension $T$ de la perche et la réaction normale de la piste.
\item Au démarrage, l'accélération vaut $0{,}5\mss$. Calculer la tension.
\item À $3\ms$, le skieur lâche la perche. Calculer sa décélération (TCI), la distance qu'il parcourt encore en montant et la durée correspondante.
\item Retrouver cette distance avec le théorème de l'énergie cinétique.
\item À vitesse constante, sur $200\U{m}$ de piste : calculer les travaux de la tension, du poids et des frottements. Vérifier le TEC.
\end{enumerate}

\newpage\sujet{Évaluation 2 du chapitre (120 min)}

\Exo{28}{Aller et retour sur un plan incliné}{7}
Un solide de $0{,}5\U{kg}$ est lancé depuis le bas $O$ d'un plan incliné de $30^{\circ}$, vers le haut, à $6\ms$. Une force de frottement constante $f=0{,}5\U{N}$ s'oppose au mouvement. On prend $g=10\mss$.


\begin{enumerate}[label=\textbf{\arabic*.},leftmargin=*,start=1]\item Appliquer le TCI pendant la montée : calculer l'accélération et la réaction normale.
\item Calculer la distance parcourue avant l'arrêt et la durée de la montée.
\item Montrer que le solide redescend, puis calculer son accélération pendant la descente.
\item Calculer sa vitesse en repassant en $O$, puis la retrouver par le TEC sur l'aller-retour.
\item Calculer l'énergie mécanique perdue sur l'aller-retour. Que devient-elle ?
\end{enumerate}

\Exo{29}{Le choc des billes}{6}
Une bille A de $200\U{g}$ roulant à $4\ms$ heurte une bille B de $600\U{g}$ immobile, sur un plan horizontal. Le choc est élastique.


\begin{enumerate}[label=\textbf{\arabic*.},leftmargin=*,start=1]\item Calculer la vitesse du centre d'inertie du système {A + B} avant le choc. Que vaut-elle après ? Justifier.
\item Calculer les vitesses de A et de B juste après le choc. Vérifier la vitesse de $G$.
\item Vérifier la conservation de l'énergie cinétique.
\item B glisse ensuite sur un sol où s'exerce une force de frottement constante de $0{,}3\U{N}$. Quelle distance parcourt-elle avant de s'arrêter ?
\item Si les billes étaient restées collées, quelle aurait été leur vitesse ? Quelle énergie aurait été perdue ?
\end{enumerate}

\Exo{30}{Le remorquage}{7}
Une voiture de $1000\U{kg}$ tire une remorque de $500\U{kg}$ par une barre horizontale, sur une route droite et horizontale. Les frottements valent $200\U{N}$ sur la voiture et $100\U{N}$ sur la remorque. On prend $g=10\mss$.


\begin{enumerate}[label=\textbf{\arabic*.},leftmargin=*,start=1]\item Le convoi roule à vitesse constante. Calculer la tension de la barre et la force motrice.
\item Au démarrage, l'accélération vaut $1{,}2\mss$. Calculer la tension de la barre (système : la remorque) et la force motrice (système : le convoi).
\item Partie du repos, quelle vitesse le convoi atteint-il après $50\U{m}$ ? Vérifier avec le TEC appliqué au convoi.
\item Le centre d'inertie de la remorque est à $6\U{m}$ derrière celui de la voiture. Situer le centre d'inertie du convoi et calculer sa quantité de mouvement à $20\ms$.
\item À $20\ms$ constante, la barre casse. Calculer la décélération de la remorque et la distance qu'elle parcourt avant de s'arrêter.
\end{enumerate}

\newpage
\mntitreserie{}{Corrigé de la série P2}

\sujet{Partie A · Centre d'inertie et référentiels}

\Exo{1}{Le système Terre-Lune}{}
Origine au centre de la Terre : $x_G=\dfrac{81M_L\times0+M_L\times384\,000}{82\,M_L}\approx4683\U{km}$. C'est moins que $6400\U{km}$ : $G$ est \textbf{à l'intérieur de la Terre}.


\Exo{2}{Comparer des quantités de mouvement}{}
\begin{enumerate}[label=\textbf{\arabic*.},leftmargin=*,start=1]\item $p=1000\times20=2\times10^4\U{kg\cdot m\cdot s^{-1}}$.
\item $v=\dfrac{2\times10^4}{10^4}=2\ms=7{,}2\U{km/h}$.
\end{enumerate}

\Exo{3}{Choisir le référentiel}{}
a) Géocentrique. b) Héliocentrique. c) Terrestre.


\Exo{4}{L'ascenseur}{}
Forces : poids et tension. Mouvement rectiligne uniforme : $\vec T+\vec P=\vec 0$, donc $T=mg=4000\U{N}$.


\Exo{5}{L'haltère (synthèse)}{10}
\begin{enumerate}[label=\textbf{\arabic*.},leftmargin=*,start=1]\item Origine au centre de la boule 1 : $x_G=\dfrac{2\times0+1\times25+3\times50}{6}=\dfrac{175}{6}\approx29{,}2\U{cm}$. \textbf{(5 pts)}
\item $P=Mv_G=6\times0{,}5=3\U{kg\cdot m\cdot s^{-1}}$. \textbf{(3 pts)}
\item Les boules n'ont pas la même masse : $G$ est décalé vers la plus lourde, la boule 2. \textbf{(2 pts)}
\end{enumerate}

\Exo{6}{Le bloc tiré à vitesse constante (synthèse)}{10}
\begin{enumerate}[label=\textbf{\arabic*.},leftmargin=*,start=1]\item Dans un référentiel galiléen, si les forces se compensent, le centre d'inertie est au repos ou en MRU (et réciproquement). \textbf{(2 pts)}
\item Forces : $\vec P$, $\vec F$, $\vec R_N$, $\vec f$. MRU : leur somme est nulle. Horizontale : $f=F=12\U{N}$. Verticale : $R_N=mg=40\U{N}$. \textbf{(4 pts)}
\item $R=\sqrt{40^2+12^2}\approx41{,}8\U{N}$ ; $\tan\varphi=\dfrac{12}{40}$, donc $\varphi\approx16{,}7^{\circ}$. \textbf{(4 pts)}
\end{enumerate}

\sujet{Partie B · Le théorème du centre d'inertie}

\Exo{7}{L'ascenseur qui démarre}{}
Axe vertical vers le haut : $T-mg=ma$, donc $T=m(g+a)=500\times12=6000\U{N}$. C'est plus que le poids ($5000\U{N}$) : il faut une force supplémentaire pour accélérer.


\Exo{8}{Sans frottement}{}
$a=\dfrac{F}{m}=6\mss$ ; $v=12\ms$ ; $d=\frac12\times6\times4=12\U{m}$.


\Exo{9}{La montée}{}
Axe vers le haut du plan : $a=-g\sin\alpha=-5\mss$. $0-36=2\times(-5)\times d$, donc $d=3{,}6\U{m}$. $0=-5t+6$, donc $t=1{,}2\U{s}$.


\Exo{10}{Mesurer un frottement}{}
$f=m(g\sin\alpha-a)=0{,}5\times(3{,}42-2{,}4)\approx0{,}51\U{N}$.


\Exo{11}{Le plan incliné puis le plan horizontal (synthèse)}{10}
\begin{enumerate}[label=\textbf{\arabic*.},leftmargin=*,start=1]\item Forces : $\vec P$, $\vec R_N$, $\vec f$. Sur $(Ox)$ : $mg\sin\alpha-f=ma_1$, donc $a_1=5-2=3\mss$. Sur $(Oy)$ : $R_N=mg\cos\alpha\approx8{,}66\U{N}$. \textbf{(4 pts)}
\item $v_B=\sqrt{2\times3\times4}=\sqrt{24}\approx4{,}90\ms$ ; $t=\dfrac{v_B}{a_1}\approx1{,}63\U{s}$. \textbf{(3 pts)}
\item $-f=ma_2$, donc $a_2=-2\mss$. $0-24=2\times(-2)\times d$, donc $d=6\U{m}$. \textbf{(3 pts)}
\end{enumerate}

\Exo{12}{La luge tirée par une corde (synthèse)}{10}
\begin{enumerate}[label=\textbf{\arabic*.},leftmargin=*,start=1]\item Forces : $\vec P$, $\vec R_N$, $\vec T$, $\vec f$. TCI : $\vec P+\vec R_N+\vec T+\vec f=m\vec a$. \textbf{(2 pts)}
\item $T\cos30^{\circ}-f=ma$, donc $a=\dfrac{86{,}6-50}{20}\approx1{,}83\mss$. \textbf{(3 pts)}
\item $R_N+T\sin30^{\circ}-mg=0$, donc $R_N=200-50=150\U{N}$. La corde tire aussi un peu vers le haut et soulage le sol. \textbf{(3 pts)}
\item $v=1{,}83\times5\approx9{,}15\ms$. \textbf{(2 pts)}
\end{enumerate}

\sujet{Partie C · Les théorèmes énergétiques}

\Exo{13}{La distance de freinage}{}
\begin{enumerate}[label=\textbf{\arabic*.},leftmargin=*,start=1]\item TEC : $0-\frac12\times1000\times400=W(\vec f)$, donc $W(\vec f)=-2\times10^5\U{J}$ ($\vec P$ et $\vec R$ ne travaillent pas).
\item $-f\,d=-2\times10^5$, donc $d=25\U{m}$.
\end{enumerate}

\Exo{14}{La chute libre}{}
TEC : $\frac12mv^2=mgh$, donc $v=\sqrt{2gh}=\sqrt{400}=20\ms$.


\Exo{15}{Le pendule simple}{}
La masse descend de $h=\ell(1-\cos60^{\circ})=0{,}5\U{m}$. Conservation de $E_m$ : $\frac12mv^2=mgh$, donc $v=\sqrt{2\times10\times0{,}5}=\sqrt{10}\approx3{,}16\ms$.


\Exo{16}{Le toboggan}{}
$E_m$ en haut : $0{,}5\times10\times1{,}8=9\U{J}$ ; en bas : $\frac12\times0{,}5\times25=6{,}25\U{J}$. $E_m$ a diminué : il y a des frottements ; énergie perdue : $2{,}75\U{J}$.


\Exo{17}{Le skieur (synthèse)}{10}
\begin{enumerate}[label=\textbf{\arabic*.},leftmargin=*,start=1]\item $W(\vec P)=mgh=80\times10\times10=8000\U{J}$ ; $W(\vec R)=0$ (réaction perpendiculaire au déplacement). \textbf{(3 pts)}
\item $\frac12mv_B^2=8000$, donc $v_B=\sqrt{200}\approx14{,}1\ms$. \textbf{(3 pts)}
\item Sur $BC$ : $0-8000=-200\,d$, donc $d=40\U{m}$. \textbf{(4 pts)}
\end{enumerate}

\Exo{18}{La balle et l'énergie mécanique (synthèse)}{10}
\begin{enumerate}[label=\textbf{\arabic*.},leftmargin=*,start=1]\item $E_m=\frac12\times0{,}2\times64+0{,}2\times10\times1{,}5=6{,}4+3=9{,}4\U{J}$. \textbf{(3 pts)}
\item Au sommet : $mgz_{\max}=9{,}4$, donc $z_{\max}=4{,}7\U{m}$. Au sol : $\frac12\times0{,}2\,v^2=9{,}4$, donc $v=\sqrt{94}\approx9{,}70\ms$. \textbf{(4 pts)}
\item Au sommet réel : $E_m=0{,}2\times10\times4{,}2=8{,}4\U{J}$ ; perte : $1\U{J}$. Montée de $2{,}7\U{m}$ : $f=\dfrac{1}{2{,}7}\approx0{,}37\U{N}$. \textbf{(3 pts)}
\end{enumerate}

\sujet{Partie D · La conservation de la quantité de mouvement}

\Exo{19}{Le canon}{}
$v=-\dfrac{5\times200}{500}=-2\ms$ : le canon recule à $2\ms$.


\Exo{20}{Sauter d'une barque}{}
$70\times2+140\,v=0$, donc $v=-1\ms$ : la barque part à $1\ms$ dans le sens opposé.


\Exo{21}{Masses égales}{}
$v_1'=\dfrac{m-m}{2m}v_1=0$ et $v_2'=\dfrac{2m}{2m}v_1=v_1$ : la première s'arrête, la seconde repart avec la vitesse $v_1$.


\Exo{22}{La balle dans le bloc}{}
\begin{enumerate}[label=\textbf{\arabic*.},leftmargin=*,start=1]\item $0{,}02\times300=1\times v'$, donc $v'=6\ms$.
\item Avant : $\frac12\times0{,}02\times9\times10^4=900\U{J}$ ; après : $\frac12\times1\times36=18\U{J}$. Perte : $882\U{J}$, soit $98$ \% de l'énergie.
\end{enumerate}

\Exo{23}{Le patineur et le ballon (synthèse)}{10}
\begin{enumerate}[label=\textbf{\arabic*.},leftmargin=*,start=1]\item Les forces extérieures sont les poids et la réaction de la glace, verticales, qui se compensent (frottements négligeables pendant le lancer). \textbf{(2 pts)}
\item $\vec P=\vec 0$ avant : $50\,v+2\times10=0$, donc $v=-0{,}4\ms$ (sens opposé au ballon). \textbf{(4 pts)}
\item TEC : $0-\frac12\times50\times0{,}16=-f\times2$, donc $f=2\U{N}$. \textbf{(4 pts)}
\end{enumerate}

\Exo{24}{Les deux pendules (synthèse)}{10}
\begin{enumerate}[label=\textbf{\arabic*.},leftmargin=*,start=1]\item Conservation de $E_m$ : $v_1=\sqrt{2gh}=\sqrt{4}=2\ms$. \textbf{(2 pts)}
\item Masses égales : $v_1'=0$ et $v_2'=2\ms$. $B_2$ remonte à $h=\dfrac{v_2'^2}{2g}=20\U{cm}$. \textbf{(4 pts)}
\item $0{,}1\times2=0{,}2\,v'$, donc $v'=1\ms$ ; $h'=\dfrac{1}{20}=5\U{cm}$. $E_c$ : $0{,}2\U{J}$ avant, $\frac12\times0{,}2\times1=0{,}1\U{J}$ après : $0{,}1\U{J}$ perdu. \textbf{(4 pts)}
\end{enumerate}

\sujet{Évaluation 1 du chapitre}

\Exo{25}{La piste ABC (questions 1 à 3 : bac 2024)}{7}
\begin{enumerate}[label=\textbf{\arabic*.},leftmargin=*,start=1]\item TEC de $A$ à $B$ : $\frac12mv_B^2=mg\ell\sin\alpha-f\ell$, donc $v_B=\sqrt{2\ell\left(g\sin\alpha-\dfrac fm\right)}$. \textbf{(1,5 pt)}
\item TEC de $A$ à $C$ : $0=mg\ell\sin\alpha-2f\ell$, donc $f=\dfrac{mg\sin\alpha}{2}=0{,}5\U{N}$. \textbf{(1,5 pt)}
\item Sur $AB$ : $a=g\sin\alpha-\dfrac fm=2{,}5\mss$ ; $\ell=\frac12at_B^2$, donc $t_B=\sqrt{0{,}8}\approx0{,}89\U{s}$ ($v_B=\sqrt5\approx2{,}24\ms$). \textbf{(1 pt)}
\item Sur $BC$ : $-f=ma_2$, $a_2=-2{,}5\mss$ ; durée $\dfrac{v_B}{2{,}5}\approx0{,}89\U{s}$ (le solide parcourt bien $\dfrac{5}{5}=1\U{m}$). \textbf{(1 pt)}
\item $R_N=mg\cos\alpha\approx1{,}73\U{N}$ ; $R=\sqrt{1{,}73^2+0{,}5^2}\approx1{,}80\U{N}$. \textbf{(1 pt)}
\item $E_m(A)=mg\ell\sin\alpha=1\U{J}$ ; $E_m(C)=0$. Perte $1\U{J}$ $=f\times2\ell$ : c'est le travail des frottements (changé en chaleur). \textbf{(1 pt)}
\end{enumerate}

\Exo{26}{Le pendule balistique}{6}
\begin{enumerate}[label=\textbf{\arabic*.},leftmargin=*,start=1]\item Conservation de $E_m$ après le choc : $v'=\sqrt{2gh}=1\ms$. \textbf{(1,25 pt)}
\item Balle et bloc restent ensemble : choc \textbf{mou}. $\vec P$ conservée : $0{,}01\,v=1\times1$, donc $v=100\ms$. \textbf{(1,5 pt)}
\item Avant : $\frac12\times0{,}01\times10^4=50\U{J}$ ; après : $\frac12\times1\times1=0{,}5\U{J}$. Perte $49{,}5\U{J}$, soit $99$ \%. \textbf{(1,25 pt)}
\item {fusil + balle} pseudo-isolé : $2\,v_f+0{,}01\times100=0$, donc $v_f=-0{,}5\ms$. \textbf{(1 pt)}
\item $v_G=\dfrac{P}{M}=\dfrac{1}{1}=1\ms$ avant et après : le système est isolé pendant le choc, $\vec P=M\vec v_G$ ne change pas. \textbf{(1 pt)}
\end{enumerate}

\Exo{27}{Le téléski}{7}
\begin{enumerate}[label=\textbf{\arabic*.},leftmargin=*,start=1]\item Principe de l'inertie (MRU) : $T=mg\sin\alpha+f\approx239{,}4+50=289{,}4\U{N}$ ; $R_N=mg\cos\alpha\approx658\U{N}$. \textbf{(1,5 pt)}
\item $T-mg\sin\alpha-f=ma$, donc $T\approx289{,}4+35=324{,}4\U{N}$. \textbf{(1 pt)}
\item $a=-\left(g\sin\alpha+\dfrac fm\right)\approx-4{,}13\mss$ ; $d=\dfrac{9}{2\times4{,}13}\approx1{,}09\U{m}$ ; $t=\dfrac{3}{4{,}13}\approx0{,}73\U{s}$. \textbf{(1,5 pt)}
\item $0-\frac12\times70\times9=-(mg\sin\alpha+f)\,d$, donc $d=\dfrac{315}{289{,}4}\approx1{,}09\U{m}$. \textbf{(1,5 pt)}
\item $W(\vec T)=289{,}4\times200\approx57\,880\U{J}$ ; $W(\vec P)=-mg\times200\sin\alpha\approx-47\,880\U{J}$ ; $W(\vec f)=-10\,000\U{J}$ ; somme nulle $=\Delta E_c$. \textbf{(1,5 pt)}
\end{enumerate}

\sujet{Évaluation 2 du chapitre}

\Exo{28}{Aller et retour sur un plan incliné}{7}
\begin{enumerate}[label=\textbf{\arabic*.},leftmargin=*,start=1]\item Axe vers le haut : $-mg\sin\alpha-f=ma_1$, $a_1=-(5+1)=-6\mss$ ; $R_N=mg\cos\alpha\approx4{,}33\U{N}$. \textbf{(1,5 pt)}
\item $d=\dfrac{36}{12}=3\U{m}$ ; $t=1\U{s}$. \textbf{(1,5 pt)}
\item $mg\sin\alpha=2{,}5\U{N}>f$ : il redescend. Axe vers le bas : $a_2=5-1=4\mss$. \textbf{(1 pt)}
\item $v^2=2\times4\times3=24$, $v\approx4{,}90\ms$. TEC : $\frac12\times0{,}5\,v^2-9=-0{,}5\times6$, donc $v^2=24$. \textbf{(1,5 pt)}
\item $E_m$ : $9\U{J}$ au départ, $6\U{J}$ au retour ; $3\U{J}$ perdus, changés en chaleur par les frottements. \textbf{(1,5 pt)}
\end{enumerate}

\Exo{29}{Le choc des billes}{6}
\begin{enumerate}[label=\textbf{\arabic*.},leftmargin=*,start=1]\item $v_G=\dfrac{0{,}2\times4}{0{,}8}=1\ms$ ; le système est isolé pendant le choc : $\vec P$, donc $\vec v_G$, ne change pas. \textbf{(1 pt)}
\item $v_A'=\dfrac{0{,}2-0{,}6}{0{,}8}\times4=-2\ms$ ; $v_B'=\dfrac{0{,}4}{0{,}8}\times4=2\ms$ ; $v_G=\dfrac{-0{,}4+1{,}2}{0{,}8}=1\ms$. \textbf{(1,5 pt)}
\item Avant : $1{,}6\U{J}$ ; après : $0{,}4+1{,}2=1{,}6\U{J}$. \textbf{(1 pt)}
\item TEC : $0-1{,}2=-0{,}3\,d$, donc $d=4\U{m}$. \textbf{(1,25 pt)}
\item $v'=v_G=1\ms$ ; $E_c=\frac12\times0{,}8\times1=0{,}4\U{J}$ : $1{,}2\U{J}$ perdus. \textbf{(1,25 pt)}
\end{enumerate}

\Exo{30}{Le remorquage}{7}
\begin{enumerate}[label=\textbf{\arabic*.},leftmargin=*,start=1]\item Remorque en MRU : $T=100\U{N}$ ; convoi en MRU : $F=200+100=300\U{N}$. \textbf{(1,5 pt)}
\item Remorque : $T-100=500\times1{,}2$, $T=700\U{N}$. Convoi : $F-300=1500\times1{,}2$, $F=2100\U{N}$. \textbf{(1,5 pt)}
\item $v=\sqrt{2\times1{,}2\times50}=\sqrt{120}\approx11\ms$. TEC : $2100\times50-300\times50=90\,000\U{J}=\frac12\times1500\times120$. \textbf{(1,5 pt)}
\item Origine sur $G_{\text{voiture}}$, axe vers l'arrière : $x_G=\dfrac{500\times6}{1500}=2\U{m}$ derrière la voiture. $P=1500\times20=3\times10^4\U{kg\cdot m\cdot s^{-1}}$. \textbf{(1,25 pt)}
\item $-100=500\,a$, $a=-0{,}2\mss$ ; $d=\dfrac{400}{0{,}4}=1000\U{m}$. \textbf{(1,25 pt)}
\end{enumerate}
\label{mnlast}
\end{document}
